% === FIGURE LEGENDS ===
% Listed separately from the figure floats (MRM style: captions go in a list at
% the end of the manuscript text, not inside the figure files). Each number is
% auto-synced to its float via \ref to the figure's \label, so the list stays
% correct even if the figures are reordered.
%TC:ignore
\section*{Figure Legends}

\noindent\textbf{Figure~\ref{fig:fisher}.}
Fisher-information analysis of the eight-bin spectral discretization.
(a)~Correlation matrix $C_{jk} = F_{jk}/\sqrt{F_{jj}F_{kk}}$ of the Fisher
information $\mathbf{F} = \Umat^T\Umat/\sigma^2$ (Eq.~\ref{eq:fisher}); off-diagonal
entries near $\pm 1$ mark pairs of components the b-value sampling cannot separate well.
(b)~Per-component estimation uncertainty at the cohort-median SNR $= 303$
(log scale): the unconstrained Cram\'er--Rao lower bound (light gray), the
Bayesian (van Trees) CRLB under a Gaussian approximation of the half-normal
prior ($\sigma_R = 3.16$; dark gray), and the empirical NUTS posterior standard
deviation (orange), the last sampled under the true half-normal prior with
non-negativity. Annotated factors mark the prior gain
(gray, unconstrained\,$\to$\,Bayesian) and the constraint gain
(orange, Bayesian\,$\to$\,NUTS).
(c)~Individual component decay curves $R_j \exp(-bD_j)/K$ at equal fractions,
with horizontal noise floors at SNR $= 50$, $100$, and $303$.

\medskip\noindent\textbf{Figure~\ref{fig:validation}.}
Validation of the spectral inference on simulated ground truth across spectrum
shape and noise level.
Six ground-truth spectra (rows A--F: cohort-mean normal and tumor, a uniform
spectrum, a concentrated $\delta$ at $D = 0.75\;\umsmm$, a reversed-tumor
spectrum, and a bimodal spectrum) are each reconstructed at SNR $= 50$, $300$,
and $1000$ (columns; SNR $300$ (@Stephan: In Fig 1 you say it is 303 (You could also change Fig 1 to 300).)is the cohort median).
In each panel the bars are the ground truth (grey), the tuned-MAP estimate
($\lambda = 10^{-3}$, green), and the NUTS posterior mean (orange); error bars
are $\pm 1$ standard deviation across 100 noise realizations (those with
split-$\hat{R} > 1.05$ excluded, $\geq 81$ retained per condition; the per-panel
$\hat{R}_{\max}$ annotated).
The bottom row shows noise recovery, $\hat{\sigma}/\sigma_{\mathrm{true}}$
(dashed reference $= 1$), for the jointly inferred NUTS $\hat{\sigma}$ and the
MAP residual plug-in (off-scale values shown with an arrow).

\medskip\noindent\textbf{Figure~\ref{fig:spectra}.}
Cohort diffusivity spectra by tissue type and anatomical zone.
For each diffusivity bin, box plots show the cohort distribution of the per-ROI
estimated water fraction $R_j$ (box: interquartile range; line: median;
whiskers: $1.5\times$ IQR), with MAP (green) alongside NUTS posterior-mean
estimates (orange) at the tuned regularization $\lambda = 10^{-3}$.
Top row: normal; bottom row: tumor. Left column: peripheral zone (PZ); right:
transition zone (TZ). Per-ROI spectra are shown in the supplementary atlas
(Figure~\ref{fig:atlas}).

\medskip\noindent\textbf{Figure~\ref{fig:roc}.}
Leave-one-out cross-validated ROC for tumor-versus-normal detection in the
peripheral zone (PZ, left; $n = 81$) and transition zone (TZ, right; $n = 68$).
The four thick reference curves---ADC (black) and the spectral classifier on
three bin subsets: all eight fractions (orange solid), the two outer fractions (@Stephan: I would remove "the" and just generically call them outer fractions with the specification in the parenthesis. (The problem is that 'the' outer fraction would mean 0.25 and 20, not 0.25 and 3.)
$\{D = 0.25,\,3.0\;\umsmm\}$ (orange dashed), and the five inner fractions with
the two outer (@Stephan: "all outer" (0.25, 3,0, 20.0)) compartments removed (orange dotted)---pass through the identical
NUTS-feature leave-one-out logistic-regression pipeline ($C = 1.0$).
The thin curves are raw single-feature ROCs (tumor-positive orientation): the
outer components $D = 0.25\;\umsmm$ (restricted, teal (@Stephan: teal is green-blue. I would call this "turquoise" or "light-blue". Non-English people will anyway not be familiar with the term teal.)) and $D = 3.0\;\umsmm$
(free water, brown); $D = 2.0\;\umsmm$ (goldenrod (@Stephan: dark yellow)), $D = 1.5\;\umsmm$ (indigo (@Stephan: blue (consider swapping with 3.0)),
and $D = 20\;\umsmm$ (magenta); and the remaining intermediate bins as a faint
grey bundle.
AUCs with 95\% confidence intervals are reported in Table~\ref{tab:auc}.

\medskip\noindent\textbf{Figure~\ref{fig:adc_discriminant}.}
ADC versus the spectral discriminant score (the logistic-regression linear
predictor from the eight spectral fractions) across ROIs, by anatomical zone
(columns: PZ, $n = 81$, left; TZ, $n = 68$, right) and spectral estimator
(rows: NUTS posterior-mean, top; constrained-MAP at tuned $\lambda = 10^{-3}$,
bottom).
All panels share common axes; the discriminant is fit (@Stephan: explain abbreviation "OLS" that is shown in Figure) separately per zone and
estimator, so its scale is comparable only within a panel (@Stephan: Looks to me as the vertical scale is the same for all panels.), whereas ADC
($x$-axis) is the same physical quantity throughout.
Tumor ROIs are red and normal ROIs blue; the black line is the
ordinary-least-squares fit, and per-panel Pearson $r$ with a 95\% bootstrap
confidence interval is annotated (all $|r| > 0.95$).

\medskip\noindent\textbf{Figure~\ref{fig:sensitivity}.}
Per-bin structure of the spectral discriminant and the ADC sensitivity vector,
by anatomical zone (PZ, TZ).
Bars give the standardized tumor-versus-normal logistic-regression coefficient
per diffusivity bin---the change in tumor log-odds per standard deviation of
that fraction, so heights are comparable across bins of differing scale
(Methods)---using NUTS features, $C = 1.0$, the same fit as
Figure~\ref{fig:adc_discriminant}, normalized to $\max|w| = 1$, with 95\%
bootstrap confidence intervals. (@Stephan: Do you or Sandy understand this? I pass. @Patrick: First I think this should maybe be in main text, considering that we had said to use figure descriptoins primarily for descritpive purposes right, second maybe there is a simpler way we can explain the standardization here ... I think we already have something in the meain text, coordinate here so that people undersatnd)
Charcoal (@Stephan: black) diamonds give the ADC sensitivity $-\partial\text{ADC}/\partial R_j$
at the average tumor spectrum (normalized, sign-aligned to the tumor
direction).
Each bar is colored by the bin's within-ROI posterior CV (light
$=$ well identified to dark $=$ poorly identified), a per-ROI identifiability
distinct from the cohort-level coefficient stability shown by the bootstrap
intervals.

\medskip\noindent\textbf{Figure~\ref{fig:spectrum_ggg}.}
NUTS posterior-mean diffusivity spectrum averaged across ROIs by Gleason Grade
grouping against normal-tissue baseline (dark gray,
$n = 109$). Left: GGG\,$=\,1$ (green, $n = 8$) versus GGG\,$\geq 2$ (red,
$n = 21$). Right: GGG\,$=\,2$ (orange, $n = 12$) versus GGG\,$\geq 3$ (purple,
$n = 9$). Shaded bands are 95\% confidence intervals of the group mean
(Student's $t$, $\text{df} = n - 1$).
Eleven of the 40 tumor ROIs (7 with GGG\,$=\,0$ and 4 ungraded) are not shown.

\medskip\noindent\textbf{Figure~\ref{fig:uncertainty}.}
Per-ROI classifier predictions in the peripheral zone (top, $n = 81$) and
transition zone (bottom, $n = 68$), each sorted by predicted probability.
The leave-one-out classifier (the same NUTS-feature pipeline as shown in
Figure~\ref{fig:roc}, $C = 1$) is applied to every NUTS posterior draw of the
eight fractions, yielding a posterior distribution over $P(\mathrm{tumor})$:
markers at the posterior-mean probability (circles normal, squares tumor) with
error bars indicating the propagated 90\% credible interval; correctly classified ROIs in
pale colors, misclassified in saturated colors.

\medskip\noindent\textbf{Figure~\ref{fig:pixelwise}.}
Pixel-wise spectral decomposition of a single peripheral-zone prostate DWI
slice.
(A)~Anatomical reference ($b = 0$).
(B,\,C)~NUTS posterior-mean restricted-diffusion ($D = 0.25\;\umsmm$) and
free-water ($D = 3.0\;\umsmm$) fractions.
(D)~Conventional ADC map (restricted diffusion appears dark).
(E,\,F)~Per-voxel discriminant score for the two-bin
$\{D = 0.25,\,3.0\;\umsmm\}$ and full eight-bin peripheral-zone classifiers
(same classification as applied in Figure~\ref{fig:roc}), centered on the decision boundary.
(H,\,I)~Posterior standard deviation of the score (500 Monte Carlo draws) for
the two- and eight-bin detectors on a shared scale.
%TC:endignore
